% SPDX-License-Identifier: Apache-2.0 WITH LLVM-exception
% "int considered harmful"
%
% Build:  pdflatex -interaction=nonstopmode int-considered-harmful.tex   (x3)
% Run from the docs/ directory; listings are read from int-considered-harmful/.

\documentclass[11pt,a4paper]{article}

\usepackage[T1]{fontenc}
\usepackage[utf8]{inputenc}
\usepackage{lmodern}
\usepackage{microtype}
\usepackage[margin=2.5cm,top=2.6cm,bottom=2.6cm]{geometry}
\usepackage{xcolor}
\usepackage{listings}
\usepackage{booktabs}
\usepackage{array}
\usepackage{amsmath}
\usepackage{amssymb}
\usepackage{enumitem}
\usepackage{fancyhdr}
\usepackage{titlesec}
\usepackage{caption}
\usepackage{longtable}
\usepackage[hidelinks,colorlinks=true,linkcolor=linkblue,urlcolor=linkblue,
            citecolor=linkblue]{hyperref}

% ---------------------------------------------------------------- colours ---
\definecolor{linkblue}{HTML}{1A4F8A}
\definecolor{codebg}{HTML}{F6F7F9}
\definecolor{codeframe}{HTML}{D8DCE2}
\definecolor{termbg}{HTML}{F2F0EC}
\definecolor{termframe}{HTML}{D5CFC4}
\definecolor{kw}{HTML}{0B5394}
\definecolor{cmt}{HTML}{6B7A86}
\definecolor{str}{HTML}{A03030}
\definecolor{num}{HTML}{2E6F40}
\definecolor{rulegray}{HTML}{888888}
\definecolor{badred}{HTML}{9B1C1C}
\definecolor{goodgreen}{HTML}{1F6B3B}

% ---------------------------------------------------------------- headers ---
\pagestyle{fancy}
\fancyhf{}
\renewcommand{\headrulewidth}{0.3pt}
\fancyhead[L]{\small\itshape int considered harmful}
\fancyhead[R]{\small\thepage}

\titleformat{\section}{\Large\bfseries}{\thesection}{0.7em}{}
\titleformat{\subsection}{\large\bfseries}{\thesubsection}{0.6em}{}
\titlespacing*{\section}{0pt}{1.6em}{0.7em}
\titlespacing*{\subsection}{0pt}{1.2em}{0.5em}

\setlist[itemize]{leftmargin=1.3em,topsep=0.3em,itemsep=0.25em}
\setlist[enumerate]{leftmargin=1.6em,topsep=0.3em,itemsep=0.25em}

% --------------------------------------------------------------- listings ---
\lstdefinestyle{cpp}{
  language=C++,
  basicstyle=\ttfamily\footnotesize,
  keywordstyle=\color{kw}\bfseries,
  commentstyle=\color{cmt}\itshape,
  stringstyle=\color{str},
  numberstyle=\tiny\color{rulegray},
  backgroundcolor=\color{codebg},
  frame=single,
  rulecolor=\color{codeframe},
  framesep=5pt,
  xleftmargin=6pt,
  xrightmargin=2pt,
  breaklines=true,
  breakatwhitespace=true,
  showstringspaces=false,
  columns=fullflexible,
  keepspaces=true,
  tabsize=2,
  aboveskip=0.9em,
  belowskip=0.9em,
  morekeywords={constexpr,consteval,static_assert,noexcept,auto,using,
                nullptr,uint8_t,int16_t,int64_t,uint16_t,size_t,concept,
                requires,co_await,decltype},
}

\lstdefinestyle{term}{
  basicstyle=\ttfamily\footnotesize,
  backgroundcolor=\color{termbg},
  frame=single,
  rulecolor=\color{termframe},
  framesep=5pt,
  xleftmargin=6pt,
  xrightmargin=2pt,
  breaklines=true,
  showstringspaces=false,
  columns=fullflexible,
  keepspaces=true,
  aboveskip=0.7em,
  belowskip=0.9em,
}

\lstset{style=cpp}

\newcommand{\code}[1]{\texttt{\small #1}}
\newcommand{\file}[1]{\texttt{\small #1}}

% caption for a code/output pair
\newcommand{\exhibit}[1]{\vspace{0.2em}\noindent{\small\itshape #1}\par\vspace{-0.4em}}

% ------------------------------------------------------------ bench tables ---
\newenvironment{benchtable}[1]{%
  \par\vspace{0.6em}\noindent{\small\bfseries #1}\par\vspace{0.2em}
  \small
  \begin{tabular}{r r r l}
  \toprule
  rel. & ns/op & ins/op & implementation \\
  \midrule
}{%
  \bottomrule
  \end{tabular}\par\vspace{0.8em}
}

% variant carrying the "does this row enforce the range?" column
\newenvironment{benchtablec}[1]{%
  \par\vspace{0.6em}\noindent{\small\bfseries #1}\par\vspace{0.2em}
  \small
  \begin{tabular}{r r r c l}
  \toprule
  rel. & ns/op & ins/op & enforces? & implementation \\
  \midrule
}{%
  \bottomrule
  \end{tabular}\par\vspace{0.8em}
}

% ================================================================= title ====
\title{\vspace{-1.5cm}\textbf{\code{int} considered harmful}\\[0.4em]
       \large A case against the default integer, and a type that fixes it}
\author{Peter Neiss}
\date{26 July 2026}

\begin{document}
\maketitle
\thispagestyle{plain}

\begin{abstract}
\noindent
C++'s \code{int} is not a number. It is a platform-dependent bit pattern whose
behaviour is undefined at both ends of its range, which silently changes width
and signedness in the middle of expressions, which rounds without being asked,
and which carries no record of what it is counting. Every one of those
properties is a defect, and all of them are load-bearing in code that ships
today.

This paper does two things. First it puts \code{int} in the dock: eleven
hazards, each a complete program, compiled and run, with the real output
reproduced. One of them causes GCC to emit \emph{no machine code at all} for
\code{main}. It then makes the same case, more briefly, against reaching for
\code{double} instead --- the reflex fix that trades one kind of silent
wrongness for another.

Second, it presents \code{inside}, a header-only C++23 library in which the
range and the step size of a number live in its type. Arithmetic widens its
result type at compile time, so overflow is not merely detected but
unreachable; division by zero is a value you can test rather than a signal that
kills the process; fractional values are exact rationals rather than binary
approximations; and what happens at the edge of a range is a policy you state
in the declaration. On performance we find no single answer: against code
providing the same guarantee, \code{inside} is by turns \emph{faster} (it
discharges at compile time checks that hand-written code must perform at
runtime), exactly free (byte-identical machine code), and substantially slower
(exact rational arithmetic, reproducible transcendentals). We report all three,
with the programs that produced them.
\end{abstract}

\vspace{0.5em}
\noindent\rule{\textwidth}{0.4pt}
\vspace{0.3em}

{\small
\noindent\textbf{A note on evidence.} Every program in this paper lives in
\file{docs/int-considered-harmful/} in the \code{inside} repository. Every
output block was captured by running the program, not written by hand. All
measurements come from the repository's generated \file{docs/performance.md} or
from named commits; none are estimates. Compiler is GCC 15.2.0 on
\mbox{x86-64} Linux unless stated otherwise. \code{inside} is alpha software and
its API may change.

\smallskip
\noindent\textbf{On AI assistance.} This paper was drafted with
\href{https://claude.com/claude-code}{Claude Code}, which wrote the example
programs, compiled and ran them to capture the outputs reproduced here,
extracted the benchmark figures from the repository's generated reports, and
produced the typeset document. The \code{inside} library itself was likewise
developed with Claude Code. The argument, the framing and the final review are
the author's; the claims were checked against the compiler and the repository
rather than taken on the model's word, and the reader is invited to do the
same --- every program is included so that they can.
}

\vspace{0.8em}
\noindent\rule{\textwidth}{0.4pt}

\tableofcontents
\newpage

% ============================================================================
\section{The thesis}

Dijkstra's 1968 letter did not argue that \code{goto} produced wrong answers.
It argued that \code{goto} destroyed the correspondence between the program on
the page and the process it described, so that a reader could no longer reason
about what was true at a given point. The objection was about
\emph{tractability}, not arithmetic.

The case against \code{int} is narrower and, in a sense, worse. \code{int} does
produce wrong answers. It produces them silently, it produces them in code that
compiles without a warning at maximum diagnostic settings, and --- the part
practitioners consistently underestimate --- it produces them by causing the
compiler to \emph{delete code the programmer wrote}, including the code written
to check for the very failure that occurs.

The root of it is a mismatch of models. The programmer thinks of \code{int} as
$\mathbb{Z}$, the integers. The standard defines it as a finite set with
undefined behaviour at the boundary, and permits the optimiser to assume the
boundary is never reached. Those two models agree over the range where testing
usually happens and diverge exactly where it usually does not.

What follows from the mismatch is a family of defects that no amount of
discipline fully removes, because the information needed to check the operation
--- what range is this value allowed to occupy? what is it counting? what
should happen if it leaves? --- is not present in the program at the point
where the operation happens. It exists only in the programmer's head and, if
you are fortunate, in a comment.

\code{inside}'s proposal is to put that information in the type, where the
compiler can act on it. The remainder of this paper establishes that the
problem is real (\S\ref{sec:indictment}--\S\ref{sec:double}), that the usual
defences do not solve it (\S\ref{sec:mitigations}), and that a type-level
answer does (\S\ref{sec:inside}--\S\ref{sec:correctness}), at a cost that is
sometimes nothing at all and sometimes substantial
(\S\ref{sec:performance}).

% ============================================================================
\section{The indictment}\label{sec:indictment}

\subsection{Signed overflow is undefined, and the optimiser believes it}

The first thing to establish is that signed overflow is not wraparound. It is
undefined behaviour, which means the compiler may assume it does not occur and
optimise on that assumption. Consider a function that asks whether adding one
to a number makes it bigger.

\lstinputlisting[firstline=13,lastline=18]{int-considered-harmful/h01_overflow_ub.cpp}

\noindent
Since \code{x + 1 > x} can only be false if \code{x + 1} overflows, and
overflow is undefined, the compiler concludes the expression is always true. It
does not warn. It deletes the comparison:

\begin{lstlisting}[style=term]
$ g++ -std=c++23 -O0 h01_overflow_ub.cpp -o h01 && objdump -d h01
0000000000001149 <grows(int)>:
    1149:   endbr64
    114d:   push   %rbp
    114e:   mov    %rsp,%rbp
    1151:   mov    %edi,-0x4(%rbp)     <- argument stored...
    1154:   mov    $0x1,%eax           <- ...and ignored. Return 1.
    1159:   pop    %rbp
    115a:   ret
\end{lstlisting}

\noindent
Note the optimisation level: \textbf{\code{-O0}}. This is not an aggressive
release-build transformation that a debug build would protect you from. GCC
folds the comparison during translation, before optimisation proper begins. The
check cannot be written in a way that survives, because the check is the thing
being deleted.

\begin{lstlisting}[style=term]
INT_MAX            = 2147483647
grows(INT_MAX)     = yes, always
INT_MAX + 1        = -2147483648
\end{lstlisting}

\noindent
The last two lines are worth reading together. The program reports that adding
one to \code{INT\_MAX} always produces a larger value, and then prints the
smaller value it produced.

\subsection{The same assumption deletes an entire function}

That was a folded comparison. The assumption compounds. Here is a loop whose
induction variable starts near the top of the range:

\lstinputlisting[firstline=12,lastline=26]{int-considered-harmful/h02_overflow_loop.cpp}

\noindent
At \code{-O0} the counter wraps to a negative value, \code{i > 0} goes false,
and the loop exits after three iterations. At \code{-O2} the compiler reasons
that \code{i > 0} can never become false without overflow, and overflow is
undefined, therefore the loop is infinite, therefore the code after it is
unreachable, therefore \code{main} has no reachable exit. What it emits is:

\begin{lstlisting}[style=term]
$ g++ -std=c++23 -O2 -Wall h02_overflow_loop.cpp -o h02
$ objdump -d h02

0000000000001040 <main>:
    1040:   endbr64
    1044:   cs nopw 0x0(%rax,%rax,1)
    104e:   xchg   %ax,%ax

0000000000001050 <_start>:
\end{lstlisting}

\noindent
That is the whole of \code{main}: an endbranch landing pad and five bytes of
alignment padding. There is no loop, no \code{printf}, and \textbf{no
\code{ret} instruction}. Execution falls out of \code{main} into whatever
follows it in the text segment:

\begin{lstlisting}[style=term]
$ ./h02
Segmentation fault (core dumped)          [exit 139]

$ ./h02_compiled_at_O0
iterations = 3
\end{lstlisting}

\noindent
\code{-Wall} produced no diagnostic. The same source, the same compiler, two
optimisation levels: one prints \code{3}, the other crashes before executing a
single instruction of the program. This is the strongest available argument
that undefined behaviour is not a theoretical concern about edge-case values.
It is a compile-time contract, and breaking it does not corrupt your result ---
it deletes your program.

\subsection{Two's complement is asymmetric}

There are $2^{31}$ negative \code{int} values and $2^{31}-1$ positive ones.
Negation therefore has no correct answer at one point in the range, and neither
does absolute value.

\lstinputlisting[firstline=14,lastline=19]{int-considered-harmful/h03_intmin.cpp}
\begin{lstlisting}[style=term]
INT_MIN            = -2147483648
INT_MAX            = 2147483647
-INT_MIN           = -2147483648   <- still negative
abs(INT_MIN)       = -2147483648   <- still negative
\end{lstlisting}

\noindent
\code{std::abs} is a function whose entire purpose is to return a non-negative
value, and it returns a negative one. It is not a bug in the standard library;
there is no value it could return that would be correct.

\subsection{Division has two undefined inputs and no error value}

Floating point has \code{NaN} and infinities. Integer division has neither, so
its two invalid inputs --- division by zero, and \code{INT\_MIN / -1}, whose
true answer is \code{INT\_MAX + 1} --- cannot be represented as results. On
x86-64 both raise \code{SIGFPE} and terminate the process.

\begin{lstlisting}[style=term]
$ ./h04_divzero 0
about to evaluate 10 / 0 ...
Floating point exception (core dumped)          [exit 136]

$ ./h04_divzero 1
about to evaluate -2147483648 / -1 ...
Floating point exception (core dumped)          [exit 136]
\end{lstlisting}

\noindent
The second case is the one that surprises people: a division whose operands are
both perfectly ordinary-looking, with no zero in sight, that kills the process.
Any code that divides by a runtime value must check for two special cases, and
almost none does.

\subsection{Mixed signed/unsigned comparison}

When an \code{int} meets an \code{unsigned} in a comparison, the usual
arithmetic conversions make both unsigned. A negative value becomes an enormous
positive one.

\lstinputlisting[firstline=13,lastline=24]{int-considered-harmful/h05_signed_unsigned.cpp}
\begin{lstlisting}[style=term]
-1 < 0u            = false
v.size()           = 0
v.size() - 1       = 18446744073709551615   <- not -1
loop body entered  = 4 times on an EMPTY vector
\end{lstlisting}

\noindent
The loop iterates over an empty container. Without the escape hatch it would
run roughly $1.8\times10^{19}$ times, indexing far past the end of the
allocation on the first pass. This is among the most common sources of
exploitable memory-safety bugs in C++ code, and it arises from a
\emph{comparison}.

\subsection{Small types do not stay small}

Integral promotion means the operands of an expression are frequently not the
types you wrote, and the result is frequently not the type of the operands.

\lstinputlisting[firstline=12,lastline=28]{int-considered-harmful/h06_promotion.cpp}
\begin{lstlisting}[style=term]
uint8 * uint8 -> int? yes
200 * 200          = 40000
stored back in u8  = 64   <- truncated
char is signed?    yes (platform-dependent)
sizeof(int)        = 4 bytes here; the standard promises >= 2
\end{lstlisting}

\noindent
Three separate portability hazards sit in that output. The multiplication of
two 8-bit values happens in \code{int}, so the intermediate is fine but the
store truncates. The signedness of bare \code{char} is
implementation-defined --- code that treats it as either is wrong somewhere.
And \code{sizeof(int)} is 4 here but the standard guarantees only 2, which is
not academic: it is 2 on a great deal of shipping embedded hardware.

\subsection{Division truncates, and not in the direction you expect}

\lstinputlisting[firstline=10,lastline=17]{int-considered-harmful/h07_truncdiv.cpp}
\begin{lstlisting}[style=term]
 7 / 2             = 3   (exact answer 3.5)
-7 / 2             = -3   (rounds toward zero, not down)
 7 % 2             = 1
-7 % 2             = -1   <- negative remainder
mean(90,90,91)     = 90   (true mean 90.333...)
\end{lstlisting}

\noindent
Integer division rounds toward zero, which is neither floor nor
round-to-nearest. It is the only rounding mode available, it is applied
silently, and it is the wrong one for most arithmetic people actually want:
computing an average, splitting a bill, scaling a value by a percentage.
The remainder is negative for negative operands, which breaks the common
idiom \code{if (x \% n == k)} for indexing into cyclic structures.

\subsection{Shifts}

\lstinputlisting[firstline=9,lastline=21]{int-considered-harmful/h08_shift.cpp}
\begin{lstlisting}[style=term]
1 << 31            = -2147483648   <- undefined behaviour
-8 >> 1            = -4
\end{lstlisting}

\noindent
Shifting a bit into the sign position of a signed type is undefined; so is
shifting by an amount greater than or equal to the width, and shifting by a
negative amount. The value printed above happens to be what the hardware does.
It is not what the standard promises, and the optimiser is not obliged to agree
with the hardware.

\subsection{Narrowing is silent}

\lstinputlisting[firstline=11,lastline=24]{int-considered-harmful/h11_narrowing.cpp}
\begin{lstlisting}[style=term]
velocity (64-bit)  = 40000
stored as 16-bit   = -25536   <- sign flipped
int 300 -> char    = 44
\end{lstlisting}

\noindent
A positive velocity becomes a negative one. This is, in miniature, the Ariane 5
Flight 501 failure of June 1996: a 64-bit floating-point horizontal-velocity
value converted to a 16-bit signed integer, in inertial-reference software
carried over unchanged from Ariane 4, where the value's range had been provably
smaller. The conversion overflowed 37 seconds after ignition, the guidance
computers shut down, and the vehicle was destroyed with its payload. The
software was correct for the range it had been written against; nothing in the
type system recorded what that range was.

\subsection{The midpoint bug}

\lstinputlisting[firstline=9,lastline=20]{int-considered-harmful/h10_binsearch.cpp}
\begin{lstlisting}[style=term]
lo                 = 2147483644
hi                 = 2147483646
(lo + hi) / 2      = -3   <- outside [lo, hi]
lo + (hi - lo) / 2 = 2147483645   <- correct
\end{lstlisting}

\noindent
The midpoint of two values inside the range lands outside it. This is the bug
Jon Bentley described in \emph{Programming Pearls} (1986) and Joshua Bloch
rediscovered in the JDK in 2006, where \code{Arrays.binarySearch} had carried it
for nine years. If a bug of this shape can live that long in one of the most
scrutinised standard libraries in existence, written by people who knew about
it, the problem is not insufficient care.

\subsection{The semantic hazard: \code{int} does not know what it counts}

Every hazard so far has been about range. This one is about meaning, and it is
the one no sanitizer will ever catch, because the arithmetic is valid.

\lstinputlisting[firstline=8,lastline=27]{int-considered-harmful/h09_units.cpp}
\begin{lstlisting}[style=term]
subtotal (cents)   = 3683
shipping (dollars) = 5
total              = 3688   <- charged 5 cents for shipping
compiles clean, runs clean, bills the customer wrong.
\end{lstlisting}

\noindent
An \code{int} holding cents and an \code{int} holding dollars are the same
type. Adding them is well-defined, warning-free, sanitizer-clean, and wrong.
The Mars Climate Orbiter was lost in 1999 to precisely this class of error
across a software interface --- pound-force-seconds where newton-seconds were
expected. No range check would have caught it, because no value was out of
range. The information that would have caught it --- what the number
\emph{means} --- was never written down anywhere the compiler could see.

% ============================================================================
\section{``Just use \code{double}''}\label{sec:double}

The common response to the preceding section is to reach for floating point.
\code{double} has a much wider range, it does not trap on division by zero, and
it represents fractions. It is a reasonable instinct and it is the right tool
for a genuine class of problems. But it is not a fix for the problems above; it
substitutes a different family of silent wrongness, and this section is
deliberately shorter than the last one only because the failures are more
widely known --- not because they are less serious.

\subsection{Decimal fractions are not representable}

\begin{lstlisting}[style=term]
0.1 + 0.2          = 0.30000000000000004441
0.3                = 0.29999999999999998890
0.1 + 0.2 == 0.3   ? false
\end{lstlisting}

\noindent
Neither \code{0.1} nor \code{0.2} nor \code{0.3} exists in binary floating
point. The literals are the nearest representable neighbours, and the sum of
the first two is not the neighbour of the third.

\subsection{The error accumulates}

\lstinputlisting[firstline=7,lastline=16]{int-considered-harmful/d02_money_drift.cpp}
\begin{lstlisting}[style=term]
10000 x $0.01      = 100.00000000001425348728
expected           = 100.00000000000000000000
exactly $100.00    ? false
error              = 0.00000000001425348728
\end{lstlisting}

\noindent
Ten thousand transactions --- a quiet morning for a payment system --- and the
ledger no longer balances to the cent. The error is small, which is precisely
what makes it dangerous: it survives testing and shows up in reconciliation.

\subsection{Addition is not associative}

\lstinputlisting[firstline=8,lastline=17]{int-considered-harmful/d03_assoc.cpp}
\begin{lstlisting}[style=term]
(a + b) + c        = 1.0
a + (b + c)        = 0.0
equal              ? false
\end{lstlisting}

\noindent
The same three values, the same operator, two groupings, two answers. The
practical consequence is that the result of a summation depends on the order
the data was visited in. Parallelise a reduction across a different number of
threads, or let the optimiser re-associate a loop, and the answer changes.
Results stop being reproducible --- not across platforms, but across
\emph{runs}.

\subsection{Integers stop counting at $2^{53}$}

\lstinputlisting[firstline=8,lastline=19]{int-considered-harmful/d04_precision.cpp}
\begin{lstlisting}[style=term]
2^53               = 9007199254740992.0
2^53 + 1           = 9007199254740992.0   <- unchanged
2^53 + 1 == 2^53   ? true
(1e16 + 1) - 1e16  = 0.0   (should be 1)
\end{lstlisting}

\noindent
Beyond $2^{53}$ a \code{double} cannot represent consecutive integers, so a
counter silently stops incrementing --- a real hazard for identifiers and
timestamps in systems that pass numbers through JSON, where \code{double} is
the only numeric type. The second line shows catastrophic cancellation:
subtracting near-equal quantities annihilates every digit that mattered.

\subsection{\code{NaN} breaks equality and ordering}

\begin{lstlisting}[style=term]
0.0 / 0.0          = -nan
1.0 / 0.0          = inf   <- no error reported
nan == nan         ? false   <- breaks equality
nan < 1            ? false
nan > 1            ? false   <- neither; breaks sorting
nan + 1000         = -nan   <- contaminates the whole total
\end{lstlisting}

\noindent
\code{double} does not trap, which is often described as an advantage. The cost
is that an invalid operation produces a value that propagates silently through
every subsequent computation, and that violates the strict-weak-ordering
requirement \code{std::sort} relies on --- sorting a range containing
\code{NaN} is undefined behaviour, and can walk off the end of the buffer.

\medskip
\noindent
The honest summary is that floating point is excellent at what it was designed
for: approximating continuous quantities across an enormous dynamic range. It
is a poor fit for money, for counting, for identifiers, and for anything
requiring reproducibility. Reaching for it to escape \code{int}'s problems
means adopting a different set.

% ============================================================================
\section{Why the usual defences fall short}\label{sec:mitigations}

Before proposing something new, it is worth being precise about why the
existing tools do not close the gap.

\paragraph{Sanitizers are runtime instruments.} UBSan genuinely detects signed
overflow, and it is valuable. But it detects it only when the offending input
actually occurs, in a build that carries the instrumentation, and it cannot
detect what the optimiser already deleted. Running our very first example under
UBSan:

\begin{lstlisting}[style=term]
$ g++ -std=c++23 -O2 -fsanitize=undefined h01_overflow_ub.cpp && ./a.out
h01_overflow_ub.cpp:23:14: runtime error: signed integer overflow:
    2147483647 + 1 cannot be represented in type 'int'
INT_MAX            = 2147483647
grows(INT_MAX)     = yes, always
INT_MAX + 1        = -2147483648
\end{lstlisting}

\noindent
UBSan reports line 23 --- the \code{printf} argument. It says nothing about
\code{grows()}, where the comparison was folded away at translation time. There
was no overflow left to instrument, because there was no arithmetic left. The
diagnostic misses the defect that motivated the example.

\paragraph{\code{-ftrapv} and \code{-fwrapv} change the language.}
\code{-fwrapv} defines overflow as wrapping, which removes the undefined
behaviour and the optimiser's licence with it. That is a real improvement, and
it makes the wrong answer reliable rather than unpredictable: the midpoint bug
still computes a negative midpoint. It is a fix for the UB, not for the
arithmetic.

\paragraph{Assertions are manual, and unenforced.} A precondition check is only
present where somebody remembered to write it, expresses a range the type
system does not know, and is frequently compiled out in release builds. It is
documentation with a runtime cost, not a guarantee.

\paragraph{Static analysis is partial by construction.} Range analysis across
function boundaries and through data structures is undecidable in general.
Analysers are useful and they are not complete; they cannot tell you that this
\code{int} is cents and that one is dollars, because nothing in the program
says so.

\paragraph{Safe-integer libraries solve one hazard.} Wrapper types that check
for overflow on every operation (\code{boost::safe\_numerics},
\code{SafeInt}) genuinely fix hazards \S2.1--\S2.2. They do not address units,
they do not address rounding policy, they do not give you a range narrower than
the underlying type, and they detect at runtime what could have been rejected
at compile time.

\medskip
\noindent
The common thread: each tool operates at a point where the necessary
information has already been discarded. The range this value may occupy, the
quantity it measures, the rounding rule that applies, and what should happen at
the boundary are all decisions the programmer made and then failed to record in
a form the compiler can act upon. The proposal in the next section is simply
to record them.

% ============================================================================
\section{\code{inside}: putting the range in the type}\label{sec:inside}

\code{inside} is a header-only C++23 library built
around a single template:

\begin{lstlisting}
template<grid G, policy_flag P>
struct inside;                        // include/beman/inside/core.hpp:62
\end{lstlisting}

\noindent
The \code{grid} is a compile-time value carrying an interval (lower and upper
inside) and a \emph{notch} --- the step size, held as an exact rational. Every
value of the type is exactly $\text{Lower} + k \cdot \text{Notch}$ for some
integer $k$. The \code{policy} says what happens when a value would leave the
range. Storage is selected automatically: the narrowest of
\code{uint8}\ldots\code{int64}, \code{float}/\code{double}, or an exact
fraction, that can represent every point on the grid.

\begin{lstlisting}
using namespace beman::inside;

using pct    = inside<{0, 100}>;                    // integers 0..100
using money  = inside<{{0, 1'000'000}, notch<1,100>}, round_nearest>;  // cents
using angle  = inside<{0, 359}, wrap>;              // modular
using sample = inside<{{-1, 1}, notch<1,16384>}, round_nearest>;       // Q1.14
\end{lstlisting}

\noindent
Three properties follow, and between them they account for most of
\S\ref{sec:indictment}.

\subsection{Arithmetic widens, so overflow is unreachable}

The result type of an operation is computed at compile time from the operand
grids. Adding two values in $[10, 255]$ yields a type whose range is
$[20, 510]$ --- wide enough that no representable input can overflow it. There
is no overflow to be undefined, no check to perform, and nothing for the
optimiser to exploit.

\lstinputlisting[firstline=13,lastline=34]{int-considered-harmful/b01_no_overflow.cpp}
\begin{lstlisting}[style=term]
a + b              = 236
a - b              = -204   (no wraparound)
(lo + hi) / 2      = 2147483645   (still in range)
\end{lstlisting}

\noindent
Three things happened there. The sum was verified at compile time by
\code{static\_assert}, so a wrong answer would have been a build failure rather
than a test failure. The difference of two \emph{unsigned-ranged} values
produced $-204$: because the result grid is $[-245, 245]$, the type widened to
a signed representation rather than wrapping. And the midpoint computation from
\S2.10 --- the one that produced $-3$ in \code{int} --- stayed in range,
because \code{lo + hi} was evaluated in a grid wide enough to hold it before
the division.

\subsection{Division by zero is a value, not a signal}

Where the divisor's grid can include zero, division returns
\code{std::expected<inside, errc>}. Where the grid provably excludes zero, it returns a
plain value with nothing to unwrap --- the check is elided because the type
proves it unnecessary.

\lstinputlisting[firstline=14,lastline=36]{int-considered-harmful/b02_divzero.cpp}
\begin{lstlisting}[style=term]
7 / 3              = 2 1/3   (exact, not 2)
div(7, 3, trunc)   = 2
10 / 0             = division by zero   (no SIGFPE)
\end{lstlisting}

\noindent
Note the first line. The default result of division is \emph{exact}: $7/3$ is
represented as the rational $2\tfrac{1}{3}$, not truncated to $2$. Truncation
is still available, but you ask for it by name, per call. The failure that
terminated the process in \S2.4 is now a value the caller can test.

\subsection{Rounding is stated, not inherited}

\lstinputlisting[firstline=16,lastline=34]{int-considered-harmful/b07_rounding.cpp}
\begin{lstlisting}[style=term]
native  -7 / 2     = -3   (toward zero, always)
round_to_nearest   = -4
floored            = -4
ceiled             = -3
truncated          = -3
exact              = -3.5
\end{lstlisting}

\noindent
Each mode is honoured, including for negative operands, where the distinction
between truncation and flooring actually bites. The repository pins this with a
property test asserting $(a/b) \cdot b + a \bmod b = a$ across whole grids for
every mode, including negative divisors
(\file{tests/test\_div\_round.cpp:129--152}).

\subsection{The edge is a policy you declare}

Out-of-range values are only reachable when assigning into a narrower type ---
never from arithmetic. What happens then is part of the declaration:

\lstinputlisting[firstline=18,lastline=46]{int-considered-harmful/b03_policies.cpp}
\begin{lstlisting}[style=term]
checked  200 ->    value outside interval
clamp    150 ->    100
wrap     370 ->    10
try_make  10 ->    value outside interval
with_clamp() 150 -> 100
\end{lstlisting}

\noindent
The wrapping behaviour that \code{int} gives you unasked and undefined is
available here as \code{wrap} --- explicitly requested, well-defined, and
producing $370 \rightarrow 10$ on a $[0,359]$ grid rather than
$\text{INT\_MAX}+1 \rightarrow \text{INT\_MIN}$. The default is \code{checked},
which reports rather than storing a value it cannot represent.

\subsection{Negation reflects the type}

The asymmetry of \S2.3 disappears because negating a
$[\text{Lower},\text{Upper}]$ value yields a value of type
$[-\text{Upper},-\text{Lower}]$:

\lstinputlisting[firstline=13,lastline=28]{int-considered-harmful/b10_negation.cpp}
\begin{lstlisting}[style=term]
coldest            = -40
-coldest           = 40   (representable by construction)
-hottest           = -125
\end{lstlisting}

\noindent
There is no value whose negation falls outside the negated type, so
\code{-INT\_MIN} has no analogue. The same reasoning gives an \code{abs} that
cannot return a negative number.

\subsection{The scale is part of the type}

The units failure of \S2.11 is a failure to record what a number means.
\code{inside} records the scale in the notch, so a value on a cent grid and a
value on a dollar grid are different types, and combining them converts:

\lstinputlisting[firstline=15,lastline=33]{int-considered-harmful/b05_units.cpp}
\begin{lstlisting}[style=term]
subtotal           = $36.83
shipping           = $5
total              = $41.83   (not $36.88)
\end{lstlisting}

\noindent
The \code{int} version of this program charged five cents for shipping. Here
the dollar value is scaled onto the cent grid automatically, because the
compiler knows what each grid's step size is.

Bare integers cannot enter the arithmetic at all --- there is nothing to say
what grid they belong to --- and the diagnostic says so directly:

\begin{lstlisting}[style=term]
$ g++ -std=c++23 -I ../../include x01_bare_int_rejected.cpp
   15 |   auto bad = x + 1;      // error: no grid for the literal
include/beman/inside/arithmetic.hpp:492:27: error: static assertion failed:
   an inside can only be added to another inside: give the scalar a grid --
   write `a + 1_ins` (or `a + just<1>` / `a + one`), or `a + inside<{lo,hi}>{n}`
   for a runtime value with a known range
\end{lstlisting}

\subsection{Bounded indices}

\lstinputlisting[firstline=14,lastline=34]{int-considered-harmful/b08_index.cpp}
\begin{lstlisting}[style=term]
data[7]            = 49
index 10           = value outside interval   (no out-of-bounds access possible)
sum of all slots   = 285
\end{lstlisting}

\noindent
The bounds check happens once, at construction, rather than at every access or
not at all. An index that exists is in range by construction; one that would
not be simply fails to be created.

% ============================================================================
\section{A worked migration}\label{sec:migration}

An invoice calculation, in \code{int}. It has three defects, and the compiler
mentions none of them.

\lstinputlisting[firstline=11,lastline=23]{int-considered-harmful/m_before.cpp}
\begin{lstlisting}[style=term]
subtotal (cents)   = 268435456
tax  (8%)          = -21474836   <- overflowed
discount (1/3)     = 89478485   <- truncated
total              = 157482135

expected tax       = 21474836.48
expected discount  = 89478485.33
\end{lstlisting}

\noindent
The invoice is for \$2,684,354.56 --- large, but not implausible. Computing 8\%
of it as \code{subtotal * 8 / 100} overflows \code{int} at the multiplication,
and the tax comes out \emph{negative}. The discount silently loses its
fractional part. The total is meaningless, and every intermediate looks like a
perfectly ordinary integer.

The same calculation in \code{inside}:

\lstinputlisting[firstline=15,lastline=46]{int-considered-harmful/m_after.cpp}
\begin{lstlisting}[style=term]
subtotal           = $2684354.56
tax  (8%)          = $214748.36   (exact, no overflow)
discount (1/3)     = $894784.85   (rounded to nearest cent)
total              = $2004318.07

discount on a 1/300 grid       = $894784 64/75   (nothing was lost)
as an exact fraction of a dollar = 67108864/75
\end{lstlisting}

\noindent
What changed, and what it cost:

\begin{itemize}
\item \textbf{The overflow is gone by construction.} \code{subtotal * 0.08\_ins}
      widens to a grid that can hold the product. There is no check, because
      there is nothing to check.
\item \textbf{The tax is exact to the cent.} \code{0.08} is parsed as an exact
      rational, not as the nearest \code{double}. The product is exact, and
      assigning it into \code{money} snaps to the nearest cent because
      \code{money} was declared \code{round\_nearest}. Truncating instead would
      require saying so.
\item \textbf{The third stays exact until it is stored.} The last two lines
      prove it: stored onto a grid of thirds-of-a-cent, the discount is
      $67108864/75$ dollars exactly. The rounding in the line above it happened
      once, at the point of storage, by a stated rule --- not silently at every
      intermediate step.
\item \textbf{The friction is real but small.} Literals need a grid
      (\code{0.08\_ins}), non-representable constants need \code{frac<1,3>}
      rather than \code{1.0/3.0}, and the declaration of \code{money} is a line
      of thought that the \code{int} version never had to have. That line is
      the point: it is where the range, the scale, and the rounding rule got
      written down.
\end{itemize}

% ============================================================================
\section{Correctness}\label{sec:correctness}

\subsection{Errors caught before the program runs}

The strongest claim \code{inside} makes is not that it detects errors but that
it rejects them at compile time. Because grids are compile-time values and the
arithmetic is \code{constexpr}, whole categories of defect become build
failures.

\lstinputlisting[firstline=18,lastline=34]{int-considered-harmful/b06_compile_time.cpp}

\noindent
Every line above is a \code{static\_assert}. If any were false the file would
not build, and no test would need to run to discover it. An out-of-range
constant is likewise rejected outright:

\begin{lstlisting}[style=term]
$ g++ -std=c++23 -I ../../include x02_range_rejected.cpp
error: 'constexpr void beman::inside::policy<W,E>::report(beman::inside::errc)'
       called in a constant expression
error: call to non-'constexpr' function 'void beman::inside::detail::raise(beman::inside::errc,
       const char*)'
\end{lstlisting}

\noindent
The mechanism is worth understanding: the error path is deliberately
non-\code{constexpr}, so reaching it during constant evaluation is
automatically a compile error. The library's own test suite
(\file{tests/test\_constexpr.cpp}) is written almost entirely in
\code{STATIC\_REQUIRE} on this basis.

\subsection{The error vocabulary}

\code{inside} reports failure in one of several shapes, and which one you get is
determined by how you spell the call. The design rule is that the shape matches
the caller's role.

\begin{center}
\small
\begin{tabular}{@{}l l p{6.6cm}@{}}
\toprule
Shape & Spelling & When \\
\midrule
Plain value & \code{a + b} & The type proves the operation cannot fail. \\
\code{expected<T,errc>} & \code{a / b}, \code{try\_make}, \code{to<T>()}, \code{math::sqrt} & The operation can fail; the caller tests the result and, if it cares, reads the cause. \\
Error code & \code{x.policy(ec) = v} & Hot paths, or bulk operations where the first error should stick. \\
Throwing & default \code{checked} & Programming errors that should not be handled locally. \\
Callback & \code{on\_clamp}, \code{on\_wrap} & The event is expected and carries data (a carry, an overshoot). \\
\bottomrule
\end{tabular}
\end{center}

\noindent
The codes themselves are a small closed enumeration
(\file{detail/debug.hpp:55}):

\begin{lstlisting}
enum class errc { domain_error, division_by_zero, overflow,
                  rounding_error, not_finite };
\end{lstlisting}

\noindent
There is no \code{<system\_error>} dependency. Where exceptions are
unavailable, \code{set\_error\_handler} replaces the failure path globally,
which is what allows the library to run freestanding on bare metal.

The callback shape is more useful than it first appears, because some
``errors'' are the actual semantics of the domain. A clock is the canonical
case: seconds overflowing 59 is not a fault, it is a carry.

\begin{lstlisting}
void add_seconds(auto n)
{
  seconds.on_wrap([&](auto, auto carry) { add_minutes(carry); }) += n;
}
\end{lstlisting}

\noindent
That is the whole of the carry logic from \file{examples/clock.cpp} --- the
wrap event delivers the overshoot to the next field up, and the cascade to
minutes and hours falls out of the same pattern.

\subsection{Exactness}

Because the notch is an exact rational rather than a binary fraction, grids
like ``cents'' and ``thirds'' are represented without approximation. The money
example from \S\ref{sec:double} run in \code{inside}:

\lstinputlisting[firstline=15,lastline=30]{int-considered-harmful/b04_money.cpp}
\begin{lstlisting}[style=term]
10000 x $0.01      = $100
exactly $100.00    ? true
as a fraction      = 100/1
\end{lstlisting}

\noindent
Ten thousand additions of one cent produce exactly one hundred dollars, in four
bytes of storage --- the same computation that drifted by
$1.4\times10^{-11}$ in \code{double}. The \code{static\_assert(sizeof(money) ==
4)} in that listing is doing real work: exactness here does not require
arbitrary-precision arithmetic, because the grid is known at compile time and
the value is stored as an integer count of notches.

Floating-point storage is still available, but only where it is provably
lossless: the \code{f64} and \code{f32} flags are permitted only on
\emph{dyadic} grids --- power-of-two notch and lower bound --- so that every
grid point is exactly representable in the significand. This is enforced by
\code{static\_assert}, and a grid too fine for the requested float silently
falls back to exact rational storage rather than losing precision. The rule
that decides it is the same $2^{53}$ boundary that \S\ref{sec:double}
demonstrated.

\subsection{Reproducibility}

For integer and rational storage, \code{+ - * /} are reproducible
unconditionally --- the repository's determinism document states it plainly:
``There is no floating point on these paths, so there is nothing for the
platform to round differently.'' Conversion from \code{double} decomposes the
value from its IEEE-754 bits rather than computing with it, so it needs no FPU
and no \code{<cmath>}.

For transcendentals, \code{beman::inside::math} ships three engines behind one API, of
which the integer CORDIC engine is reproducible ``unconditionally $\ldots$ any
platform, any flags, no FPU.'' Its contract is worth quoting because it is
unusually specific:

\begin{quote}\small
Every function below produces bit-identical output $\ldots$ across compiler,
platform, optimisation level, and FP flags. 1. NO \code{<cmath>}/FPU/intrinsics;
hot paths are integer-only over int64. 2. NO runtime-derived tables $\ldots$
3. NO external code generators $\ldots$ 5. Each transcendental ships checked-in
\code{static\_assert} vectors pinning its bit-exact output.
\end{quote}

\noindent
Measured against a long-double reference, the double engine's maximum error is
exactly 0.5000 notches for essentially every function --- correctly rounded ---
and the CORDIC engine stays below 0.9 notches (\file{docs/accuracy.md}).

One honest caveat, which the library documents itself: the three engines are
reproducible individually but are \emph{not} value-identical to each other. A
golden vector captured under one engine is valid for that engine only. This is
the table-maker's dilemma, not a defect, but it does mean the engine choice is
part of your reproducibility contract.

% ============================================================================
\section{Performance}\label{sec:performance}

An abstraction that costs nothing is suspicious. This one does not have a
single cost at all, and the most important thing to say about its performance
is that \textbf{the cost is not a number, it is a range with a negative end}.
Depending on what the code is doing, using \code{inside} instead of \code{int}
can make the program faster, leave it unchanged instruction-for-instruction, or
make it substantially slower. All three regimes are measured in this section:

\begin{center}
\small
\begin{tabular}{@{}l >{\raggedright\arraybackslash}p{5.1cm} >{\raggedright\arraybackslash}p{5.9cm}@{}}
\toprule
Regime & When & Measured \\
\midrule
\textbf{Negative cost} & A chain of operations whose intermediate ranges the
compiler can discharge, replacing checks the equivalent hand-written code must
perform at runtime. &
22\% faster than native checked (\S\ref{sec:checks}); comparison at 100.2\%
of native with fewer instructions \\[0.6em]
\textbf{Zero cost} & Integer-aligned operands on the same grid, where all the
range reasoning happens in the type system. &
Byte-identical disassembly (\S\ref{sec:performance});
\code{inside<unsafe>} at 99.9\% (\S\ref{sec:fair}) \\[0.6em]
\textbf{Additional cost} & Exact rational storage, bit-reproducible
transcendentals, and element-wise checking through generic algorithms. &
\code{inside<exact>} $\approx$8$\times$ \code{double};
\code{beman::inside::math} at 9--12\% of \code{std};
\code{std::accumulate} at 2.9\% \\
\bottomrule
\end{tabular}
\end{center}

\noindent
Which regime you land in is mostly determined by how the code is written rather
than by the library: staying in inside-space lets the grid arithmetic discharge
obligations, while round-tripping through raw \code{int} at every step forces a
check at each boundary and forfeits the advantage. A single headline number ---
in either direction --- would be a misrepresentation, so this section reports
all three with the programs that produced them.

\subsection{The fast path is byte-identical to native}

The headline claim is that when both operands are integer-aligned on the same
grid, the operation lowers to a plain machine instruction. That claim is
testable, so here is the test:

\lstinputlisting[firstline=16,lastline=28]{int-considered-harmful/p01_codegen.cpp}
\begin{lstlisting}[style=term]
$ g++ -std=c++23 -O2 -I ../../include -c p01_codegen.cpp && objdump -d p01.o

0000000000000000 <inside_add(inside<...>, inside<...>)>:
   0:   endbr64
   4:   lea    (%rdi,%rsi,1),%eax
   7:   ret

0000000000000010 <native_add(int, int)>:
  10:   endbr64
  14:   lea    (%rdi,%rsi,1),%eax
  17:   ret
\end{lstlisting}

\noindent
The two functions are byte-identical. All of the range reasoning happened in
the type system and left no trace in the object code. Loops over these types
still vectorise: both the \code{inside} and the native loop emit packed SIMD
adds, the \code{inside} version on 16-bit lanes (its storage for $[0,1000]$ is
narrower than \code{int}) and the native version on 32-bit lanes.

This is not an accident that could regress unnoticed. The repository enforces
it in CI: \file{tests/check\_codegen.sh} compiles six named kernels at
\code{-O2}, disassembles them, and \textbf{fails the build if any of them
contains a \code{call} instruction}. A committed cachegrind baseline
(\file{tests/perf\_baseline.json}) guards instruction counts separately.

\subsection{What it costs, and against what}

The tables below are extracted from the repository's generated benchmark
report. \textbf{Read the \emph{rel.} column as a fraction of the first row's
throughput}: 100\% is parity with that row, and smaller is slower. Absolute
nanosecond figures are host-specific; the report itself warns that CPU
frequency scaling and turbo were active during the run, so the ratios --- not
the times --- are the signal.

The choice of baseline matters more than any individual number, and it is
easy to read these tables unfairly in \code{inside}'s favour or against it.
\textbf{The 100\% row is always plain, unchecked native \code{int}}: a bare
cast and an add, with no range enforcement of any kind. It is the fastest
thing in the table because it is not doing the job. It is the code from
\S\ref{sec:indictment} --- the code that overflows silently, that lets 300
become 44, that computes a negative midpoint. Comparing \code{inside} against
it measures the cost of correctness that the baseline never pays.

The benchmark therefore also carries two \emph{hand-written native} baselines
that do enforce a range, and these are the honest comparison
(\file{tests/bench.cpp:65}):

\begin{itemize}
\item \code{native clamped} --- the same arithmetic wrapped in
      \code{std::clamp(v, 0, 200)}.
\item \code{checked (exceptions)} --- a minimal hand-rolled bounded integer
      struct that range-checks in its constructor and in \code{operator+},
      throwing \code{std::out\_of\_range} on violation.
\end{itemize}

\noindent
Neither of those is \code{inside}; both are ordinary C++ a competent programmer
would write to get the guarantee \code{inside} provides. In the tables below the
\code{inside} rows are the ones named \code{inside<...>}, \code{beman::inside::}, or
\code{inside}; everything else is a native reference point.

\subsection{The like-for-like comparison}\label{sec:fair}

The generated report does not pair \code{inside<checked>} against an
equivalently-checked native implementation for the same operation, so the
question ``what does \code{inside} cost against code that does the same job?''
cannot be answered from it directly. \file{p02\_fair\_bench.cpp} in the
companion directory answers it by measuring all five contenders on the same
data in the same program: a summed pass over 4096 values, repeated 20{,}000
times, with inputs chosen so that no contender ever actually throws.

\begin{center}
\small
\begin{tabular}{@{}l r r c r@{}}
\toprule
Implementation & ns/op & vs.\ unchecked & enforces? & vs.\ safe native \\
\midrule
\code{native unchecked}   & 0.441 & 100.0\% & no  & --- \\
\code{native clamped}     & 0.463 &  95.4\% & yes & 100.0\% \\
\code{native checked}     & 0.519 &  85.0\% & yes &  89.2\% \\
\code{inside<checked>}     & 0.530 &  83.3\% & yes & \textbf{87.4\%} \\
\code{inside<unsafe>}      & 0.442 &  99.9\% & no  & --- \\
\bottomrule
\end{tabular}
\end{center}

\noindent
\emph{Mean of three runs; GCC 15.2.0, \code{-O2}, x86-64. The final column is
relative to \code{native clamped}, the fastest native implementation that
actually enforces the range.}

Read against the correct baseline the result is very different from what the
unchecked column suggests:

\begin{itemize}
\item \code{inside<checked>} runs at \textbf{87.4\%} of the fastest safe native
      code, and at \textbf{97.9\%} of the hand-rolled checked struct --- within
      noise of the code you would have written by hand to get the same
      guarantee.
\item \code{inside<unsafe>} runs at \textbf{99.9\%} of unchecked native, which
      is the same conclusion the disassembly in \S\ref{sec:performance}
      reached by a different route.
\item The cost of the guarantee itself, in plain C++ with no library involved,
      is 5--15\%. That is most of what the \code{inside<checked>} row is paying
      for. The library's own overhead on top of it is about 2\%.
\end{itemize}

\noindent
This is a favourable measurement pattern and it should be labelled as one: a
tight loop over contiguous data, where every contender vectorises and the
branch predictor sees the same outcome every iteration. It is not evidence
that \code{inside} is free in all circumstances --- the tables that follow show
cases where it is decidedly not. It is evidence for the narrower claim that
\emph{when you compare against code doing the same work, the gap is small},
and that most of the dramatic slowdowns reported against ``native'' are
measuring the absence of a safety check rather than the presence of a library.

\subsection{Where are the checks?}\label{sec:checks}

The table above still shows \code{inside<checked>} marginally behind the
hand-rolled native struct, which is a fair thing to be suspicious of:
\code{inside} holds strictly more compile-time information, so it should be able
to prove some checks unnecessary and delete them. It does. The reason it gains
nothing in that particular benchmark is that the benchmark gives it nothing to
prove --- every iteration converts a raw \code{int} into an inside and
immediately narrows the result back, which is the one shape where a check is
genuinely unavoidable.

Disassembling each phase separately shows exactly where a check is emitted and
where it is not:

\begin{center}
\small
\begin{tabular}{@{}l l c@{}}
\toprule
Operation & Emitted & Check? \\
\midrule
\code{bchecked(int v)}            & \code{cmp \$0xc8,\%edi; ja ...} & yes \\
\code{auto c = a + b} (kept wide) & \code{movzbl; movzbl; lea; ret} & \textbf{none} \\
\code{bchecked c = a + b}         & \code{add; cmp \$0xc8,\%rax; jg ...} & yes \\
\code{checked\_u8 c = a + b} (native) & \code{add; cmp \$0xc8,\%eax; jg ...} & yes \\
\code{bchecked b = a} (same grid)  & \code{movzbl; ret} & \textbf{none} \\
\code{wide w = a} (widening)       & \code{movzbl; ret} & \textbf{none} \\
\bottomrule
\end{tabular}
\end{center}

\noindent
A check appears in exactly two places: entering \code{inside} from a raw integer
whose range nothing knows, and narrowing a wide result into a smaller grid.
Everywhere the grid arithmetic can discharge the obligation --- same-grid
copies, widening assignments, and every intermediate in an expression --- the
check is gone, and \code{inside}'s narrowing check is instruction-for-instruction
what the native struct emits.

The consequence shows up as soon as there is a \emph{chain} of operations
rather than a single one. A native checked type must re-validate after every
step, because each intermediate has to fit back into the checked type.
\code{inside} lets the intermediates widen into grids computed at compile time
and checks only the final store. Summing four values in $[0,50]$ into a total
declared as $[0,200]$ (\file{p03\_check\_elision.cpp}):

\begin{center}
\small
\begin{tabular}{@{}l r r c@{}}
\toprule
Implementation & ns/op & vs.\ unchecked & runtime checks \\
\midrule
\code{native unchecked}  & 0.881 & 100.0\% & 0 \\
\code{native checked}    & 1.085 &  81.2\% & 3 \\
\code{inside<checked>}    & 0.888 &  \textbf{99.2\%} & \textbf{0} \\
\bottomrule
\end{tabular}
\end{center}

\noindent
\code{inside<checked>} is \textbf{22\% faster than the native checked type} here,
and within 1\% of native code that performs no checking at all --- while
providing the guarantee that the unchecked version does not. Zero checks, not
one: four values in $[0,50]$ sum to $[0,200]$, which is exactly the declared
range of the result, so the final store is proved safe too. The whole function
is:

\begin{lstlisting}[style=term]
endbr64
add    %esi,%edi
add    %edx,%edi
add    %ecx,%edi
movzbl %dil,%eax
ret
\end{lstlisting}

\noindent
That is byte-for-byte the unchecked native chain. The native checked version,
by contrast, emits three \code{cmp}/\code{jg} pairs, a cold path per throw
site, and enough register pressure to spill.

The proof is real rather than skipped: narrow the result type by one, to
$[0,199]$, and the sum no longer provably fits --- the store check reappears in
the disassembly immediately. This is the answer to ``why use a library instead
of writing the checks by hand'': the hand-written checks are the ones that
cannot be removed, because nothing in \code{int} records what has already been
established.

\begin{benchtablec}{construct (u8 [0,200])}
  100.0\% & 4.71 & 98.00 & no & \texttt{native uint8} \\
  63.5\% & 7.41 & 162.00 & yes & \texttt{native clamped} \\
  87.6\% & 5.37 & 124.00 & yes & \texttt{native checked (exceptions)} \\
  34.3\% & 13.71 & 211.00 & no & \texttt{inside<unsafe>} \\
\end{benchtablec}

\begin{benchtablec}{add (u8 [0,200])}
  100.0\% & 6.52 & 151.00 & no & \texttt{native uint8} \\
  69.2\% & 9.42 & 216.00 & yes & \texttt{native clamped} \\
  61.7\% & 10.57 & 234.00 & yes & \texttt{native checked (exceptions)} \\
  16.5\% & 39.60 & 541.00 & no & \texttt{inside<unsafe>} \\
\end{benchtablec}

\begin{benchtable}{increment (u32 [0,200000])}
  100.0\% & 2.34 & 43.00 & \texttt{native ++} \\
  24.3\% & 9.61 & 135.00 & \texttt{inside ++} \\
\end{benchtable}

\begin{benchtable}{accumulate 1000 (u32 [0,200000], per element)}
  100.0\% & 0.83 & 24.11 & \texttt{native loop} \\
  36.4\% & 2.28 & 59.12 & \texttt{inside<unsafe> loop} \\
  15.9\% & 5.20 & 101.12 & \texttt{inside<checked> loop} \\
  52.0\% & 1.59 & 34.20 & \texttt{beman::inside::sum<checked> (bulk check)} \\
  107.8\% & 0.77 & 23.12 & \texttt{std::accumulate native} \\
  2.9\% & 28.50 & 329.24 & \texttt{std::accumulate inside} \\
\end{benchtable}

\noindent
The accumulate table carries the single most actionable result in the report.
An \code{inside<checked>} element-wise loop runs at 15.9\% of native because it
range-checks every element. Using \code{beman::inside::sum}, which validates the total
instead of each addend, moves it to 52.0\% --- a 3.3$\times$ improvement from
choosing the right API for the job. The final row, \code{std::accumulate} over
\code{inside} at 2.9\%, is the cost of routing through a generic algorithm that
cannot use the bulk path.

\begin{benchtable}{comparison ==}
  100.0\% & 7.01 & 152.00 & \texttt{native int16} \\
  76.5\% & 9.17 & 186.00 & \texttt{inside same grid} \\
  100.2\% & 7.00 & 136.00 & \texttt{inside == int scalar (value raw)} \\
  100.4\% & 6.99 & 139.00 & \texttt{inside Q8.8 == int scalar (index raw)} \\
\end{benchtable}

\noindent
Comparison against a scalar reaches native parity --- and does so with
\emph{fewer} instructions than the native path (136 and 139 against 152),
because the grid tells the compiler the comparison can be done on raw storage
without a conversion.

\begin{benchtable}{f64-backed add (dyadic 1/16384 grid)}
  100.0\% & 7.13 & 147.00 & \texttt{native double} \\
  42.5\% & 16.78 & 294.00 & \texttt{inside<f64>} \\
\end{benchtable}

\begin{benchtable}{exact-backed add (rational storage)}
  100.0\% & 36.37 & 453.96 & \texttt{native fraction} \\
  505.5\% & 7.19 & 147.00 & \texttt{std double (reference)} \\
  64.8\% & 56.12 & 800.93 & \texttt{inside<exact>} \\
\end{benchtable}

\noindent
Exact rational arithmetic costs about 35\% against a hand-rolled reduced
\code{int64} fraction, and roughly 8$\times$ against \code{double} --- which is
the price of the exactness demonstrated in \S\ref{sec:correctness}, not an
inefficiency. If you do not need exact thirds, do not ask for them.

\begin{benchtable}{math: sqrt}
  100.0\% & 4.97 & 108.00 & \texttt{std double} \\
  10.5\% & 47.15 & 670.00 & \texttt{beman::inside::math inside} \\
\end{benchtable}

\begin{benchtable}{math: sin}
  100.0\% & 7.62 & 174.52 & \texttt{std double} \\
  11.8\% & 64.41 & 845.13 & \texttt{beman::inside::math inside} \\
\end{benchtable}

\noindent
The transcendentals are the most expensive thing in the library, running at
roughly 9--12\% of \code{std}. They buy bit-exact reproducibility across
platform, optimisation level and FP flags, and in the CORDIC engine's case they
run with no FPU at all. That is a real trade and it should be made
deliberately.

\begin{benchtable}{sort 10k (per element)}
  100.0\% & 281.65 & 3{,}942.63 & \texttt{native int16} \\
  80.3\% & 350.93 & 4{,}892.36 & \texttt{inside} \\
\end{benchtable}

\begin{benchtable}{transform v+1 10k (per element)}
  100.0\% & 1.33 & 38.02 & \texttt{native uint8} \\
  22.5\% & 5.94 & 102.02 & \texttt{inside u255 (16-bit storage)} \\
  23.1\% & 5.78 & 102.02 & \texttt{inside u254 (8-bit storage)} \\
\end{benchtable}

\subsection{Reading these numbers honestly}

Five conclusions, stated plainly:

\begin{enumerate}
\item \textbf{There is no single cost.} The measurements in this section span
      negative cost (22\% faster than checked native, \S\ref{sec:checks}), zero
      cost (byte-identical codegen, \S\ref{sec:performance}) and real cost
      (8$\times$ for exact rationals). Any summary of the form ``\code{inside}
      costs $N$\%'' is wrong in at least one direction; the honest question is
      always \emph{which regime does my code sit in}.
\item \textbf{Compare against code that does the same job.} The single largest
      source of misleading numbers here is the unchecked native baseline. It is
      the fastest row in every table because it enforces nothing --- it is the
      code whose failures fill \S\ref{sec:indictment}. Measured against native
      code that actually enforces the range, \code{inside<checked>} runs at
      87.4\% (\S\ref{sec:fair}), not the 16.5\% that the unchecked column in
      the imported \code{add} table implies --- and on a chain of operations it
      is 22\% \emph{faster} than the checked native type
      (\S\ref{sec:checks}).
\item \textbf{The fast paths are genuinely free.} Aligned integer arithmetic
      and scalar comparison compile to the same instructions as \code{int},
      verified by disassembly and guarded in CI. \code{inside<unsafe>} measured
      at 99.9\% of unchecked native in \S\ref{sec:fair}, by a different method,
      which is the same answer.
\item \textbf{Some paths genuinely cost, and no baseline choice rescues them.}
      Exact rational arithmetic is roughly 8$\times$ \code{double};
      \code{beman::inside::math} transcendentals run at 9--12\% of \code{std}; a
      \code{std::accumulate} over \code{inside} lands at 2.9\%. These buy
      exactness and bit-reproducibility respectively, and if you do not need
      those properties you should not pay for them.
\item \textbf{API choice dominates policy choice.} Moving from a checked
      element-wise loop to \code{beman::inside::sum} gained more (3.3$\times$) than
      switching from \code{checked} to \code{unsafe} does. Reach for the bulk
      operations before reaching for the unsafe policy.
\end{enumerate}

\noindent
The imported tables are a single run on one machine with frequency scaling
enabled; \S\ref{sec:fair} is a mean of three runs of a program included in the
companion directory. Both establish orders of magnitude, not precise ratios.
Anyone making a decision on the margin should re-run \code{perf\_report} and
\file{p02\_fair\_bench.cpp} on their own hardware.

% ============================================================================
\section{When not to use this}\label{sec:limits}

Intellectual honesty requires the other side of the ledger.

\begin{itemize}
\item \textbf{It is alpha software.} The README says so; the public API may
      change between versions. That is a real adoption risk for long-lived
      code.
\item \textbf{It requires C++23}, which rules out a good deal of existing
      work. GCC 13+, Clang 19+ (Clang 18 with libc++), MSVC 19.36+.
\item \textbf{Compile times increase.} The arithmetic is a compile-time
      computation over grids; that work has to happen somewhere, and it happens
      in your build.
\item \textbf{Error messages are template errors.} The library invests in this
      --- the \code{x01} diagnostic quoted in \S\ref{sec:inside} is a
      hand-written message naming the fix --- but a mistake off the guided path
      still produces the kind of output C++ templates produce.
\item \textbf{Widening surprises \code{auto} users.} The type of \code{a + b}
      is not the type of \code{a}. This is the mechanism that makes overflow
      unreachable, and it means \code{auto} variables in a chain of arithmetic
      acquire progressively wider types. That is correct, and it is not what
      most C++ programmers expect.
\item \textbf{Not everything needs it.} A loop counter over a small container
      is not where numeric bugs come from. The cases that justify the
      machinery are the ones where a wrong answer is expensive: money, physical
      units, fixed-point signal processing, embedded registers, control loops,
      anything crossing a system boundary.
\end{itemize}

% ============================================================================
\section{Conclusion}

The defects catalogued in \S\ref{sec:indictment} are not exotic. They are the
arithmetic behaviour of the type that C++ programmers reach for by default,
several million times a day. They are not caused by carelessness --- the
midpoint bug survived nine years in the JDK, and Ariane 5 was destroyed by code
that had been correct on the previous vehicle. They persist because the
information required to catch them is absent from the program at the point
where it would be needed.

\code{int} does not know its range, so it cannot check it. It does not know
what it counts, so it cannot object to being added to something else. It has no
rounding policy, so it truncates. It has no representation for failure, so
division by zero kills the process. And because its overflow is undefined
rather than merely wrong, the compiler is entitled to reason from the
assumption that the failure cannot occur --- which is how a program that should
print \code{3} becomes a \code{main} containing no instructions.

Writing those four facts into the type is not a large idea, and it does not
require a new language. \code{inside} demonstrates that a range, a step size and
an edge policy are enough to make overflow unreachable rather than merely
detected, to make division by zero a value, to make rounding a decision, and to
make a unit mismatch a compile error --- while emitting, on its fast paths,
byte-for-byte the same instructions as the \code{int} it replaces.

Dijkstra's closing argument was that the quality of a programmer is a
decreasing function of the density of \code{goto} statements in their
programs. That is not the claim here. The claim is narrower and easier to
check: for any numeric quantity in your program, you already know its range,
its units, and what should happen at its limits. Writing that down costs a
line. Not writing it down costs whatever the bug costs.

% ============================================================================
\appendix
\section{Building the examples}

Every program cited in this paper is in \file{docs/int-considered-harmful/} in
the repository.

\begin{lstlisting}[style=term]
# the int and double hazard programs need nothing but a compiler
g++ -std=c++23 -O2 h01_overflow_ub.cpp -o h01
g++ -std=c++23 -O0 h01_overflow_ub.cpp -o h01_O0    # compare the two

# the inside programs need the library on the include path
g++ -std=c++23 -O2 -I ../../include b01_no_overflow.cpp -o b01

# or against the committed single header, with no other -I
g++ -std=c++23 -O2 -I ../../single_include b01_no_overflow.cpp -o b01

# the two x*.cpp files are expected to FAIL to compile, by design
g++ -std=c++23 -I ../../include x01_bare_int_rejected.cpp    # error, as intended
\end{lstlisting}

\noindent
The library itself builds with CMake 3.24+:

\begin{lstlisting}[style=term]
cmake -B build && cmake --build build
ctest --test-dir build --output-on-failure
ctest --test-dir build -L example        # run the 30+ example programs
\end{lstlisting}

\section{Program index}

\begin{center}
\small
\begin{tabular}{@{}l l@{}}
\toprule
File & Demonstrates \\
\midrule
\file{h01\_overflow\_ub.cpp}      & Overflow check folded away, even at \code{-O0} \\
\file{h02\_overflow\_loop.cpp}    & \code{main} compiled to no instructions \\
\file{h03\_intmin.cpp}            & \code{-INT\_MIN} and \code{abs(INT\_MIN)} negative \\
\file{h04\_divzero.cpp}           & \code{SIGFPE} on \code{x/0} and \code{INT\_MIN/-1} \\
\file{h05\_signed\_unsigned.cpp}  & Loop over an empty container \\
\file{h06\_promotion.cpp}         & Promotion, \code{char} signedness, \code{int} width \\
\file{h07\_truncdiv.cpp}          & Truncation toward zero, negative remainder \\
\file{h08\_shift.cpp}             & Shift into the sign bit \\
\file{h09\_units.cpp}             & Cents added to dollars \\
\file{h10\_binsearch.cpp}         & Midpoint outside the range \\
\file{h11\_narrowing.cpp}         & Silent narrowing (Ariane 5) \\
\midrule
\file{d01\_decimal.cpp}           & \code{0.1 + 0.2 != 0.3} \\
\file{d02\_money\_drift.cpp}      & Accumulated drift on a running total \\
\file{d03\_assoc.cpp}             & Addition is not associative \\
\file{d04\_precision.cpp}         & $2^{53}$; catastrophic cancellation \\
\file{d05\_nan.cpp}               & \code{NaN} breaks equality and ordering \\
\midrule
\file{b01\_no\_overflow.cpp}      & Widening arithmetic; the midpoint, fixed \\
\file{b02\_divzero.cpp}           & Division as \code{expected}; exact by default \\
\file{b03\_policies.cpp}          & \code{checked} / \code{clamp} / \code{wrap} / \code{try\_make} \\
\file{b04\_money.cpp}             & Exact cents in four bytes \\
\file{b05\_units.cpp}             & Scale in the type \\
\file{b06\_compile\_time.cpp}     & Range errors as build failures \\
\file{b07\_rounding.cpp}          & Every rounding mode honoured \\
\file{b08\_index.cpp}             & Bounded indices \\
\file{b10\_negation.cpp}          & Negation reflects the grid \\
\midrule
\file{m\_before.cpp} / \file{m\_after.cpp} & The worked migration \\
\file{x01\_bare\_int\_rejected.cpp} & Bare literal rejected (expected failure) \\
\file{x02\_range\_rejected.cpp}     & Out-of-range constant rejected (expected failure) \\
\file{p01\_codegen.cpp}             & Fast path vs.\ native disassembly \\
\file{p02\_fair\_bench.cpp}         & Like-for-like benchmark vs.\ checked native \\
\file{p03\_check\_elision.cpp}      & Chain of ops: inside elides checks native cannot \\
\bottomrule
\end{tabular}
\end{center}

\section{References}

\begin{enumerate}[leftmargin=1.8em]
\item E.~W. Dijkstra, ``Go To Statement Considered Harmful,''
      \emph{Communications of the ACM}, 11(3):147--148, March 1968.
\item J.~Bentley, \emph{Programming Pearls}, Addison-Wesley, 1986 --- the
      binary search midpoint.
\item J.~Bloch, ``Extra, Extra --- Read All About It: Nearly All Binary
      Searches and Mergesorts are Broken,'' Google Research Blog, June 2006.
\item J.~L. Lions et al., \emph{ARIANE 5 Flight 501 Failure: Report by the
      Inquiry Board}, ESA/CNES, July 1996.
\item NASA, \emph{Mars Climate Orbiter Mishap Investigation Board Phase I
      Report}, November 1999.
\item D.~Goldberg, ``What Every Computer Scientist Should Know About
      Floating-Point Arithmetic,'' \emph{ACM Computing Surveys},
      23(1):5--48, 1991.
\item ISO/IEC 14882, \emph{Programming languages --- C++}, \S7.1 [expr.pre]
      (undefined behaviour on overflow), \S7.4 [conv.integral].
\item The \code{inside} repository: \file{README.md}, \file{docs/tutorial.md},
      \file{docs/policies.md}, \file{docs/arithmetic.md},
      \file{docs/determinism.md}, \file{docs/accuracy.md},
      \file{docs/performance.md}.
\end{enumerate}

\end{document}
